\documentclass{article}
\usepackage{graphicx}
\usepackage[margin=1.5cm]{geometry}
\usepackage{amsmath}

\begin{document}
\twocolumn

\title{Problem Set 2: Loops and Conditional Logic}
\author{Prof. Jordan C. Hanson}

\maketitle

\section{Exercise 1: Sorting with a For Loop}

\begin{enumerate}

\item A placement test assigns each student a math score from
0 through 5.  Students are assigned to courses according to

\[
\begin{array}{c|c}
\textbf{Score} & \textbf{Course} \\ \hline
4,5 & \text{Calculus} \\
2,3 & \text{Trigonometry} \\
0,1 & \text{Algebra}
\end{array}
\]

Begin with the following list of student tuples:

\begin{verbatim}
students = [
    ("Aaliyah", 4),
    ("Ben", 1),
    ("Camila", 3),
    ("Diego", 5),
    ("Elena", 2),
    ("Felix", 0),
    ("Grace", 3),
    ("Hector", 4),
    ("Isabel", 1),
    ("Jamal", 5),
    ("Keira", 0),
    ("Luis", 2)
]
\end{verbatim}

\begin{itemize}

\item Create the same three course structures used in the
in-class exercise:

\begin{verbatim}
algebra = ["Algebra", []]

trigonometry = [
    "Trigonometry", []
]

calculus = ["Calculus", []]
\end{verbatim}

Then combine them into the tuple

\begin{verbatim}
math_courses = (
    algebra,
    trigonometry,
    calculus
)
\end{verbatim}

\item Previously, the conditional statement had to be
repeated once for every student.  This time, use a
\texttt{for} loop to examine each tuple in
\texttt{students}:

\begin{verbatim}
for student in students:
    name = student[0]
    score = student[1]

    # conditional logic goes here
\end{verbatim}

Complete the loop using one
\texttt{if}-\texttt{elif}-\texttt{else} statement.

\item Use \texttt{append()} inside the conditional statement
to place the student's name into the correct course.
Do not manually type any student name into a course roster.

\item After the loop finishes, print

\begin{verbatim}
print(math_courses)
\end{verbatim}

and verify that all 12 students appear exactly once.

\item Print the three rosters separately using indexing:

\begin{verbatim}
print(math_courses[0])
print(math_courses[1])
print(math_courses[2])
\end{verbatim}

\end{itemize}
\end{enumerate}


\section{Exercise 2: Sorting with a While Loop}

\begin{enumerate}

\item Repeat the placement problem from Exercise 1, but this
time use a \texttt{while} loop instead of a
\texttt{for} loop.

Begin with new empty course lists:

\begin{verbatim}
algebra2 = ["Algebra", []]

trigonometry2 = [
    "Trigonometry", []
]

calculus2 = ["Calculus", []]
\end{verbatim}

and create

\begin{verbatim}
math_courses2 = (
    algebra2,
    trigonometry2,
    calculus2
)
\end{verbatim}

\begin{itemize}

\item A \texttt{while} loop requires a variable that keeps
track of which student is being processed.  Begin with

\begin{verbatim}
i = 0
\end{verbatim}

and construct a loop having the form

\begin{verbatim}
while i < 12:
    name = students[i][0]
    score = students[i][1]

    # conditional logic goes here

    i = i + 1
\end{verbatim}

\item Complete the same
\texttt{if}-\texttt{elif}-\texttt{else} decision used in
Exercise 1.

\item Explain why

\begin{verbatim}
i = i + 1
\end{verbatim}

must occur during every pass through the loop.  What would
happen if this statement were omitted?

\item Print \texttt{math\_courses2}.  The student names in
the three rosters should be identical to those obtained in
Exercise 1.

\item Compare the two approaches.  In the \texttt{for} loop,
what variable represents the current student?  In the
\texttt{while} loop, what variable determines which student
is selected from the list?

\end{itemize}
\end{enumerate}


\section{Exercise 3: Math Placement System}

\begin{enumerate}

\item Write a program called
\texttt{math\_placement.py} that collects student information
from the user and automatically constructs the three course
rosters.

The program will process exactly 10 students.  You may
assume that every score entered by the user is an integer
from 0 through 5.

\begin{itemize}

\item Begin with an empty list:

\begin{verbatim}
students = []
\end{verbatim}

Use a \texttt{while} loop and the \texttt{input()} function
to collect the name and math score of each student.

A useful starting structure is

\begin{verbatim}
i = 0

while i < 10:
    name = input("Student name: ")
    score = int(input("Math score: "))

    # store this student

    i = i + 1
\end{verbatim}

\item Store each student as a tuple of the form

\begin{verbatim}
(name, score)
\end{verbatim}

and append that tuple to the \texttt{students} list.

After the loop finishes, the structure of
\texttt{students} should be

\begin{verbatim}
[
    (name1, score1),
    (name2, score2),
    ...
    (name10, score10)
]
\end{verbatim}

\item Create the three initially empty course structures:

\begin{verbatim}
algebra = ["Algebra", []]

trigonometry = [
    "Trigonometry", []
]

calculus = ["Calculus", []]
\end{verbatim}

and combine them into

\begin{verbatim}
math_courses = (
    algebra,
    trigonometry,
    calculus
)
\end{verbatim}

\item Use a \texttt{for} loop to process the completed
\texttt{students} list.  For every student, extract the name
and score and use conditional logic to apply the placement
rules

\[
\begin{array}{c|c}
0,1 & \text{Algebra} \\
2,3 & \text{Trigonometry} \\
4,5 & \text{Calculus}.
\end{array}
\]

The conditional logic, rather than manually written student
names, must populate the three rosters.

\item After all students have been sorted, print a final
report having the general form

\begin{verbatim}
Math Placement Report

Algebra:
['student1', 'student2', ...]

Trigonometry:
['student3', 'student4', ...]

Calculus:
['student5', 'student6', ...]
\end{verbatim}

Use the course names and student lists already stored inside
\texttt{math\_courses}; do not create a second copy of the
roster information just for printing.

\item Run the program with 10 test students.  Your test data
must include at least one score from each of the three
placement ranges.

Verify that

\begin{itemize}
\item every entered student appears in exactly one course,
\item scores 0--1 are assigned to Algebra,
\item scores 2--3 are assigned to Trigonometry, and
\item scores 4--5 are assigned to Calculus.
\end{itemize}

\item Your completed program should use only concepts we
have used in class: variables, strings and integers,
\texttt{input()}, \texttt{print()}, lists, tuples,
indexing, \texttt{append()}, conditional statements, and
simple \texttt{for} and \texttt{while} loops.

Do not use dictionaries, functions, list comprehensions, or
other Python structures that have not yet been introduced.

\end{itemize}
\end{enumerate}

\end{document}